% Options for packages loaded elsewhere
\PassOptionsToPackage{unicode}{hyperref}
\PassOptionsToPackage{hyphens}{url}
\documentclass[
]{article}
\usepackage{xcolor}
\usepackage[margin=1in]{geometry}
\usepackage{amsmath,amssymb}
\setcounter{secnumdepth}{-\maxdimen} % remove section numbering
\usepackage{iftex}
\ifPDFTeX
  \usepackage[T1]{fontenc}
  \usepackage[utf8]{inputenc}
  \usepackage{textcomp} % provide euro and other symbols
\else % if luatex or xetex
  \usepackage{unicode-math} % this also loads fontspec
  \defaultfontfeatures{Scale=MatchLowercase}
  \defaultfontfeatures[\rmfamily]{Ligatures=TeX,Scale=1}
\fi
\usepackage{lmodern}
\ifPDFTeX\else
  % xetex/luatex font selection
\fi
% Use upquote if available, for straight quotes in verbatim environments
\IfFileExists{upquote.sty}{\usepackage{upquote}}{}
\IfFileExists{microtype.sty}{% use microtype if available
  \usepackage[]{microtype}
  \UseMicrotypeSet[protrusion]{basicmath} % disable protrusion for tt fonts
}{}
\makeatletter
\@ifundefined{KOMAClassName}{% if non-KOMA class
  \IfFileExists{parskip.sty}{%
    \usepackage{parskip}
  }{% else
    \setlength{\parindent}{0pt}
    \setlength{\parskip}{6pt plus 2pt minus 1pt}}
}{% if KOMA class
  \KOMAoptions{parskip=half}}
\makeatother
\usepackage{color}
\usepackage{fancyvrb}
\newcommand{\VerbBar}{|}
\newcommand{\VERB}{\Verb[commandchars=\\\{\}]}
\DefineVerbatimEnvironment{Highlighting}{Verbatim}{commandchars=\\\{\}}
% Add ',fontsize=\small' for more characters per line
\usepackage{framed}
\definecolor{shadecolor}{RGB}{248,248,248}
\newenvironment{Shaded}{\begin{snugshade}}{\end{snugshade}}
\newcommand{\AlertTok}[1]{\textcolor[rgb]{0.94,0.16,0.16}{#1}}
\newcommand{\AnnotationTok}[1]{\textcolor[rgb]{0.56,0.35,0.01}{\textbf{\textit{#1}}}}
\newcommand{\AttributeTok}[1]{\textcolor[rgb]{0.13,0.29,0.53}{#1}}
\newcommand{\BaseNTok}[1]{\textcolor[rgb]{0.00,0.00,0.81}{#1}}
\newcommand{\BuiltInTok}[1]{#1}
\newcommand{\CharTok}[1]{\textcolor[rgb]{0.31,0.60,0.02}{#1}}
\newcommand{\CommentTok}[1]{\textcolor[rgb]{0.56,0.35,0.01}{\textit{#1}}}
\newcommand{\CommentVarTok}[1]{\textcolor[rgb]{0.56,0.35,0.01}{\textbf{\textit{#1}}}}
\newcommand{\ConstantTok}[1]{\textcolor[rgb]{0.56,0.35,0.01}{#1}}
\newcommand{\ControlFlowTok}[1]{\textcolor[rgb]{0.13,0.29,0.53}{\textbf{#1}}}
\newcommand{\DataTypeTok}[1]{\textcolor[rgb]{0.13,0.29,0.53}{#1}}
\newcommand{\DecValTok}[1]{\textcolor[rgb]{0.00,0.00,0.81}{#1}}
\newcommand{\DocumentationTok}[1]{\textcolor[rgb]{0.56,0.35,0.01}{\textbf{\textit{#1}}}}
\newcommand{\ErrorTok}[1]{\textcolor[rgb]{0.64,0.00,0.00}{\textbf{#1}}}
\newcommand{\ExtensionTok}[1]{#1}
\newcommand{\FloatTok}[1]{\textcolor[rgb]{0.00,0.00,0.81}{#1}}
\newcommand{\FunctionTok}[1]{\textcolor[rgb]{0.13,0.29,0.53}{\textbf{#1}}}
\newcommand{\ImportTok}[1]{#1}
\newcommand{\InformationTok}[1]{\textcolor[rgb]{0.56,0.35,0.01}{\textbf{\textit{#1}}}}
\newcommand{\KeywordTok}[1]{\textcolor[rgb]{0.13,0.29,0.53}{\textbf{#1}}}
\newcommand{\NormalTok}[1]{#1}
\newcommand{\OperatorTok}[1]{\textcolor[rgb]{0.81,0.36,0.00}{\textbf{#1}}}
\newcommand{\OtherTok}[1]{\textcolor[rgb]{0.56,0.35,0.01}{#1}}
\newcommand{\PreprocessorTok}[1]{\textcolor[rgb]{0.56,0.35,0.01}{\textit{#1}}}
\newcommand{\RegionMarkerTok}[1]{#1}
\newcommand{\SpecialCharTok}[1]{\textcolor[rgb]{0.81,0.36,0.00}{\textbf{#1}}}
\newcommand{\SpecialStringTok}[1]{\textcolor[rgb]{0.31,0.60,0.02}{#1}}
\newcommand{\StringTok}[1]{\textcolor[rgb]{0.31,0.60,0.02}{#1}}
\newcommand{\VariableTok}[1]{\textcolor[rgb]{0.00,0.00,0.00}{#1}}
\newcommand{\VerbatimStringTok}[1]{\textcolor[rgb]{0.31,0.60,0.02}{#1}}
\newcommand{\WarningTok}[1]{\textcolor[rgb]{0.56,0.35,0.01}{\textbf{\textit{#1}}}}
\usepackage{graphicx}
\makeatletter
\newsavebox\pandoc@box
\newcommand*\pandocbounded[1]{% scales image to fit in text height/width
  \sbox\pandoc@box{#1}%
  \Gscale@div\@tempa{\textheight}{\dimexpr\ht\pandoc@box+\dp\pandoc@box\relax}%
  \Gscale@div\@tempb{\linewidth}{\wd\pandoc@box}%
  \ifdim\@tempb\p@<\@tempa\p@\let\@tempa\@tempb\fi% select the smaller of both
  \ifdim\@tempa\p@<\p@\scalebox{\@tempa}{\usebox\pandoc@box}%
  \else\usebox{\pandoc@box}%
  \fi%
}
% Set default figure placement to htbp
\def\fps@figure{htbp}
\makeatother
\setlength{\emergencystretch}{3em} % prevent overfull lines
\providecommand{\tightlist}{%
  \setlength{\itemsep}{0pt}\setlength{\parskip}{0pt}}
\usepackage{bookmark}
\IfFileExists{xurl.sty}{\usepackage{xurl}}{} % add URL line breaks if available
\urlstyle{same}
\hypersetup{
  pdftitle={Pooling Log-Odds Ratios},
  hidelinks,
  pdfcreator={LaTeX via pandoc}}

\title{Pooling Log-Odds Ratios}
\author{}
\date{\vspace{-2.5em}2026-04-17}

\begin{document}
\maketitle

\subsection{Introduction}\label{introduction}

For binary outcomes, the odds ratio (OR) is the most widely reported
effect measure across case-control studies, randomised trials with
binary endpoints, and observational comparative studies. Because the OR
is bounded below at zero with a long right tail, its sampling
distribution is markedly skewed and inverse-variance pooling on the
natural scale would be poorly behaved. The standard approach is
therefore to pool on the log scale, where the distribution is
approximately Normal, and back-transform the pooled estimate and its
confidence interval to the OR scale for reporting. Pooling on the log
scale stabilises variance, allows symmetric confidence intervals on the
analysis scale, and keeps the multiplicative-on-back-transform
interpretation that clinicians and epidemiologists expect.

\subsection{Prerequisites}\label{prerequisites}

A working understanding of odds ratios as a measure of association in
2×2 tables, inverse-variance meta-analytic weighting, and the
distinction between fixed-effect and random-effects models.

\subsection{Theory}\label{theory}

For a study with \(a_i, b_i, c_i, d_i\) in the four cells of the
exposure-by-outcome 2×2 table, the per-study log-OR and its variance are

\[\log \mathrm{OR}_i = \log \frac{a_i d_i}{b_i c_i}, \qquad \mathrm{Var}(\log \mathrm{OR}_i) = \frac{1}{a_i} + \frac{1}{b_i} + \frac{1}{c_i} + \frac{1}{d_i}.\]

A fixed-effect pool is \(\hat\theta = \sum w_i \theta_i / \sum w_i\)
with \(w_i = 1/\mathrm{Var}(\log\mathrm{OR}_i)\). A random-effects pool
adds the between-study variance \(\tau^2\) to each study's weight,
widening the confidence interval to reflect between-study uncertainty
about the underlying effect. Both pools are fit on the log scale and
back-transformed via \(\exp(\cdot)\) for reporting.

\subsection{Assumptions}\label{assumptions}

Studies are independent; the log-OR is approximately Normal (typically
holds when each cell has ≥5 events); the within-study variance is
accurately estimated; the pooled population is reasonably homogeneous
(fixed-effect) or, more commonly, is heterogeneous with effects drawn
from a Normal distribution around a common mean (random-effects).

\subsection{R Implementation}\label{r-implementation}

\begin{Shaded}
\begin{Highlighting}[]
\FunctionTok{library}\NormalTok{(metafor)}

\FunctionTok{set.seed}\NormalTok{(}\DecValTok{2026}\NormalTok{)}
\NormalTok{k }\OtherTok{\textless{}{-}} \DecValTok{10}
\NormalTok{df }\OtherTok{\textless{}{-}} \FunctionTok{data.frame}\NormalTok{(}
  \AttributeTok{ai =} \FunctionTok{rpois}\NormalTok{(k, }\DecValTok{30}\NormalTok{), }\AttributeTok{bi =} \FunctionTok{rpois}\NormalTok{(k, }\DecValTok{170}\NormalTok{),}
  \AttributeTok{ci =} \FunctionTok{rpois}\NormalTok{(k, }\DecValTok{20}\NormalTok{), }\AttributeTok{di =} \FunctionTok{rpois}\NormalTok{(k, }\DecValTok{180}\NormalTok{)}
\NormalTok{)}

\NormalTok{es }\OtherTok{\textless{}{-}} \FunctionTok{escalc}\NormalTok{(}\AttributeTok{measure =} \StringTok{"OR"}\NormalTok{, }\AttributeTok{ai =}\NormalTok{ ai, }\AttributeTok{bi =}\NormalTok{ bi,}
             \AttributeTok{ci =}\NormalTok{ ci, }\AttributeTok{di =}\NormalTok{ di, }\AttributeTok{data =}\NormalTok{ df)}

\NormalTok{fe }\OtherTok{\textless{}{-}} \FunctionTok{rma}\NormalTok{(yi, vi, }\AttributeTok{data =}\NormalTok{ es, }\AttributeTok{method =} \StringTok{"FE"}\NormalTok{)}
\NormalTok{re }\OtherTok{\textless{}{-}} \FunctionTok{rma}\NormalTok{(yi, vi, }\AttributeTok{data =}\NormalTok{ es, }\AttributeTok{method =} \StringTok{"REML"}\NormalTok{)}

\FunctionTok{cat}\NormalTok{(}\StringTok{"FE pooled OR:"}\NormalTok{, }\FunctionTok{round}\NormalTok{(}\FunctionTok{exp}\NormalTok{(fe}\SpecialCharTok{$}\NormalTok{b), }\DecValTok{2}\NormalTok{),}
    \StringTok{" 95\% CI:"}\NormalTok{, }\FunctionTok{paste}\NormalTok{(}\FunctionTok{round}\NormalTok{(}\FunctionTok{exp}\NormalTok{(}\FunctionTok{c}\NormalTok{(fe}\SpecialCharTok{$}\NormalTok{ci.lb, fe}\SpecialCharTok{$}\NormalTok{ci.ub)), }\DecValTok{2}\NormalTok{), }\AttributeTok{collapse =} \StringTok{"{-}"}\NormalTok{), }\StringTok{"}\SpecialCharTok{\textbackslash{}n}\StringTok{"}\NormalTok{)}
\FunctionTok{cat}\NormalTok{(}\StringTok{"RE pooled OR:"}\NormalTok{, }\FunctionTok{round}\NormalTok{(}\FunctionTok{exp}\NormalTok{(re}\SpecialCharTok{$}\NormalTok{b), }\DecValTok{2}\NormalTok{),}
    \StringTok{" 95\% CI:"}\NormalTok{, }\FunctionTok{paste}\NormalTok{(}\FunctionTok{round}\NormalTok{(}\FunctionTok{exp}\NormalTok{(}\FunctionTok{c}\NormalTok{(re}\SpecialCharTok{$}\NormalTok{ci.lb, re}\SpecialCharTok{$}\NormalTok{ci.ub)), }\DecValTok{2}\NormalTok{), }\AttributeTok{collapse =} \StringTok{"{-}"}\NormalTok{), }\StringTok{"}\SpecialCharTok{\textbackslash{}n}\StringTok{"}\NormalTok{)}
\end{Highlighting}
\end{Shaded}

\subsection{Output \& Results}\label{output-results}

\texttt{escalc(measure\ =\ "OR")} computes the per-study log-OR and its
variance, applying continuity correction to zero cells when needed.
\texttt{rma()} fits the meta-analysis on the log scale;
back-transformation via \texttt{exp()} yields the pooled OR and its
confidence interval on the natural scale. The random-effects model
additionally returns \(\tau^2\) and \(I^2\) describing between-study
heterogeneity.

\subsection{Interpretation}\label{interpretation}

A reporting sentence: ``Random-effects meta-analysis of 10 studies
(REML, total \(n = 4{,}000\)) gave a pooled odds ratio of 1.56 (95 \% CI
1.31 to 1.87) with modest between-study heterogeneity (\(I^2 = 22 \%\),
\(\tau^2 = 0.01\)); the fixed-effect estimate was almost identical
(1.54), confirming that heterogeneity did not materially shift the
central effect.'' Always report both the pooled estimate and the
heterogeneity metrics; readers cannot judge the pool without them.

\subsection{Practical Tips}\label{practical-tips}

\begin{itemize}
\tightlist
\item
  Always pool on the log scale; ORs are multiplicative and pooling on
  the natural scale produces biased and asymmetrically distributed
  estimates.
\item
  Apply a continuity correction (typically add 0.5 to every cell of a
  study with any zero cell) for studies with zero events;
  \texttt{metafor::escalc()} handles this automatically by default.
\item
  Choose between fixed-effect and random-effects models based on the
  substantive research question and expected sources of heterogeneity,
  not on the \(Q\)-test \(p\)-value alone --- \(I^2\) and \(\tau^2\) are
  more informative about the magnitude of heterogeneity.
\item
  Always report heterogeneity statistics (\(\tau^2\), \(I^2\), and the
  prediction interval) alongside the pooled point estimate.
\item
  Apply the Hartung--Knapp--Sidik--Jonkman adjustment for random-effects
  pooling with small \(k\) (usually fewer than 20 studies); it produces
  wider, more conservative CIs that better match coverage in
  simulations.
\item
  Mantel--Haenszel pooling (\texttt{metafor::rma.mh()}) is preferable
  for sparse-data settings with many zero cells, where continuity
  correction biases the standard inverse-variance pool.
\end{itemize}

\subsection{R Packages Used}\label{r-packages-used}

\texttt{metafor::escalc()} and \texttt{rma()} for the canonical log-OR
pooling workflow; \texttt{metafor::rma.mh()} for Mantel--Haenszel
pooling and \texttt{rma.peto()} for the Peto OR method on rare events;
\texttt{meta::metabin()} as the alternative interface with cleaner
forest-plot defaults; \texttt{metafor} Hartung--Knapp via
\texttt{test\ =\ "knha"}.

\subsection{For Reviewers}\label{for-reviewers}

\textbf{What to look for in a paper using this method.}

\begin{itemize}
\tightlist
\item
  \textbf{Common misapplications.}
\item
  \textbf{Diagnostics that should be reported but often aren't.}
\item
  \textbf{Red flags in tables and figures.}
\item
  \textbf{What to verify.}
\item
  \textbf{What an adequate Methods paragraph must contain.}
\end{itemize}

\subsection{See also --- labs in this
chapter}\label{see-also-labs-in-this-chapter}

\begin{itemize}
\tightlist
\item
  \href{labs/lab-systematic-reviews-and-prisma.Rmd}{Systematic reviews
  and PRISMA}
\item
  \href{labs/lab-meta-analysis-basics.Rmd}{Meta-analysis basics}
\item
  \href{labs/lab-network-meta-analysis.Rmd}{Network meta-analysis}
\end{itemize}

\end{document}
